\section{材料与方法}
\label{sec:methods}

\subsection{平台总体架构}

平台按"科学引擎生成结果文件、报告层消费结果文件"的原则组织（图~\ref{fig:workflow}）。数据质量控制、特征工程、受控训练、归因与统计模块之间只通过结构化结果文件（CSV/JSON/Markdown）交换信息，因此训练可以在本地或高性能计算环境中独立执行，下游分析在获得结果后离线运行。平台在分析阶段不修改任何训练产物：所有统计量与图表均由只读结果文件重新计算，随机种子、阈值与分析方案统一记录在配置文件中。工程实现细节（模块划分、接口约定、测试）见补充材料与项目仓库文档。

\subsection{案例数据集}

案例数据为 DeepCRISPR 公开数据集\citep{Chuai2018GenomeBiol}，包含四个细胞系（HCT116、HEK293T、HeLa、HL60），每条记录包含 23\,nt 靶序列、实测归一化编辑效率与四条逐位点表观遗传轨道（CTCF、DNase、H3K4me3、RRBS 的二值可及性）。过滤非法碱基后共 16\,749 条序列进入分析（表~\ref{tab:dataset}）。

数据编号经核对：全部记录的第 22、23 位均为 G，即 PAM 位于第 21--23 位（NGG），第 1--20 位为 protospacer。本文据此定义 PAM 邻近种子区为第 17--20 位、种子核心区为第 9--16 位、PAM 远端区为第 1--8 位。

\subsection{特征表示}

序列采用 23$\times$4 one-hot 编码；四条表观遗传轨道以逐位点二值形式拼接，形成 23$\times$8 输入张量（184 维展平特征）。环境组合通过\emph{通道掩码}实现：给定环境因子集合 $S$，未启用的表观通道整体置零，张量维度保持不变，从而保证不同环境组合之间的差异只反映信息可用性。Linear Regression 按统计惯例剔除 \texttt{\_T} 参照列（184$\to$161），其余模型使用完整 184 维。

\subsection{受控实验设计}

每个（模型配置 $\times$ 环境组合）在三种划分下训练：\emph{single}（单细胞系内 70\%/15\%/15\%，种子 42，四个细胞系各一次）；\emph{all}（按 held-out 细胞系分组的行）；\emph{mixed}（四细胞系混合，种子 42--45）。7 个模型配置 $\times$ 16 个环境组合 $\times$ 12 个（划分，种子）设定共产生 **1\,344 次训练运行**，全部结果在留出测试集上评估（$R^{2}$、RMSE、MAE、Pearson、Spearman）。

需要明确的数据限制：本批次中 \emph{all} 划分的 $R^{2}$ 与 \emph{single} **逐位相同**（$\max|R^{2}_{\text{all}}-R^{2}_{\text{single}}|=0$），说明预期的留一细胞系（LOCO）分支在该批次中退化，因此本文**不**将其解释为独立的跨细胞系泛化证据；真正的 LOCO 验证属于后续工作（\S\ref{sec:future}）。

\subsection{模型与归因方法}

平台使用五类模型（7 个配置）：Linear Regression（线性主效应与统计检验）、XGBoost（非线性与交互\citep{Chen2016KDD}）、MLP（一般非线性组合）、双分支 CNN（序列分支 + 环境分支，序列卷积核 $k\in\{3,5,7\}$，用于建模不同宽度的局部序列上下文）与 Transformer（跨位置交互\citep{Vaswani2017NeurIPS}）。归因方法分别为：回归系数（含 SE、$t$、$p$、FDR）、TreeSHAP\citep{Lundberg2017NeurIPS}、Integrated Gradients（IG）\citep{Sundararajan2017ICML}、CNN 的 IG 与 In Silico Mutagenesis（ISM，逐位点核苷酸替换效应）以及注意力权重。

平台强制语义边界：$\Delta R^{2}$ 表示增量预测价值；attribution 表示模型依赖强度；SNR 表示归因稳健性；$p$/FDR 仅来自具备明确零假设的统计检验（配对 bootstrap、置换检验、析因方差分析与富集检验）。**注意力权重、SNR 与树内部重要性均不作为统计显著性**，注意力在本研究中仅作支持性信息，不用于候选模式提取。

\subsection{环境因素析因分析}

对每个（划分，细胞系，模型，种子）分组定义配对增量
\begin{equation}
\Delta R^{2}_{e \mid S} = R^{2}(S \cup \{e\}) - R^{2}(S),
\end{equation}
其中 $S$ 为基线环境集合，$e$ 为新增环境因子，二者来自同一测试集。平台枚举 16 节点的 factorial lattice 的全部单因素边（同一组合可有多个父节点），节点记录 $R^{2}$/RMSE/MAE/Pearson/Spearman 均值，边记录配对的 $\Delta R^{2}$、$\Delta$MAE、$\Delta$RMSE。主效应为对全部不含 $e$ 的背景 $S$ 的条件增量均值；交互效应为 $I(a,b)=\Delta(a\mid S_0\cup\{b\})-\Delta(a\mid S_0)$ 在背景 $S_0$（不含 $a,b$）与种子上的平均。

\subsection{统计方法}

（1）\textbf{逐样本配对 bootstrap}：对每条 lattice 边，使用两侧模型在同一测试样本上的逐样本预测，以 $B=2000$ 次重采样计算 $\Delta R^{2}$ 的百分位置信区间（$\alpha=0.05$，种子 2024）；样本数不等或预测文件缺失时标记为不可用。（2）\textbf{因子级 bootstrap}：对每个环境因子，先计算各模型稳定主效应的均值（模型等权），再对该模型向量重采样，保证置信区间与点估计同口径（$n=7$ 个模型，属跨模型不确定性）。（3）\textbf{置换检验}：对逐样本损失差与条件增量使用符号翻转（sign-flip）随机化检验，$B=1000$。（4）\textbf{多重检验}：BH-FDR 仅在同一个科学问题族内校正（环境边、环境主效应、环境交互、析因 ANOVA、motif 富集各自成族）。（5）\textbf{区组析因方差分析}：以 $R^{2}$ 为响应变量，四个环境因子为二水平因子，模型、细胞系与划分作为区组因子，采用基于设计矩阵的 Type-II 边际平方和 $F$ 检验（含成对交互项），效应量报告 partial $\eta^{2}$；数据不足时返回不可用而非近似值。所有阈值统一配置（SNR 2.5/1.8/1.2；FDR 0.001/0.01/0.05；最小效应量 0.005；数值发散阈值 10.0；bootstrap $B=2000$；置换 $B=1000$），完整清单见补充表~\ref{tab:methods}。

\subsection{序列模式发现}

序列模式发现按"归因 $\rightarrow$ 片段 $\rightarrow$ 聚类 $\rightarrow$ 共识 $\rightarrow$ 富集 $\rightarrow$ 稳定性"执行。首先计算归因\emph{特异性}：位置 $p$ 上碱基 $b$ 的得分定义为其归因值减去该位置其它碱基归因均值，以避免原始幅值普遍偏高；随后按位置级分位阈值（0.90）与局部连续性要求（$\ge 2$ 个连续位置超过 0.75 分位）提取窗口（长度 4--12\,nt，每样本最多 2 个）。seqlet 按 CNN 核配置与方法分别聚类（与簇共识一致率 $\ge0.90$ 则合并，聚类后再对相似共识 $\ge0.95$ 二次合并）。共识以 IUPAC 退化码、人类可读形式与正则形式同时保存。候选模式需满足最小支持度（seqlet $\ge30$、样本 $\ge20$）。由于当前批次的归因输出为无符号幅值，模式方向不取自 attribution 符号，而由"携带该模式的序列 vs 背景序列"的实测效率对比给出，并在结果中标注来源。富集分析为可选步骤：以效率上三分位序列为前景、全部可用序列为背景，采用 Fisher 精确检验并在 motif 富集族内 BH-FDR 校正；未执行富集时不报告任何 $p$ 值。跨配置一致性通过比较不同 CNN 核与不同细胞系间共识序列的相似性（$\ge0.80$）评估；当前批次没有 Transformer IG 产物，因此 CNN 与 Transformer 的跨模型模式比较标记为不可用。

\subsection{证据整合}

平台按覆盖度、方向一致率、统计/归因强度与置信区间给出证据等级：覆盖度不足两个模型、方向冲突或置信区间跨 0 记为 Inconclusive；两个及以上模型支持、方向一致率达标且存在强统计或归因证据记为 Tier 1；其余满足覆盖度的情形为 Tier 2；仅单模型且有强证据时为 Tier 3（模型特异）。

需要强调：**证据等级反映证据的一致性、覆盖度与稳健程度，而不代表生物学效应强度**；在增量效应很小的情形下（$|\Delta R^{2}|<0.01$），高等级仅表示多个模型对该小效应的判断一致，不能解读为"强影响因素"。

\subsection{候选优先级与最小扰动方案}

在证据整合之后，平台按预先设定的标准从既有结果中筛选用于后续实验验证的候选：（i）多模型支持；（ii）效应具有方向；（iii）至少在一个人细胞系中效应较大；（iv）统计或稳健性证据非空白；（v）可设计最小扰动（单碱基替换或短序列破坏）；（vi）不依赖复杂环境操控。候选及其依据、最小扰动方案与优先级记录于表~\ref{tab:candidates}。平台在此环节只做**候选优先级排序**，不进行序列从头生成，也不声称任何候选已被实验验证。

\subsection{软件与可复现性}

平台实现为 Python 模块化分析引擎，训练与分析相互独立。所有统计量均由只读结果文件重新计算；分析方案、任务执行状态（\texttt{analysis\_status.json}）与每个任务的产物路径同时落盘；报告层与图形界面只读取这些文件。本文全部数值由 \texttt{paper/make\_assets.py} 从结果文件重新生成，并逐项记录于补充材料表~\ref{tab:claimmap} 与项目文档中的事实来源表。
